\subsection{TENSORS AND TENSOR NOTATION}

Let A be a commutative ring, M an A-module, \(\mathbf{M}^*\) the dual of M (II, § 2, no. 3) and I and J two disjoint finite sets; the A-module \(\bigotimes_{i\in \mathrm{I}\cup \mathrm{J}}\mathbf{E}_i\) , where \(\mathbf{E}_i = \mathbf{M}\) if \(i\in \mathrm{I}\) , \(\mathbf{E}_i = \mathbf{M}^*\) if \(i\in \mathrm{J}\) , is denoted by \(\mathsf{T}_{\mathsf{J}}^{\mathsf{I}}(\mathbf{M})\) ; the elements of \(\mathsf{T}_{\mathsf{J}}^{\mathsf{I}}(\mathbf{M})\) are called tensors of type (I,J) over M. They are called contravariant if \(\mathbf{J} = \varnothing\) , covariant if \(\mathbf{I} = \varnothing\) and mixed otherwise.

Let \(I'\) , \(I''\) be two subsets of I and \(J'\) , \(J''\) two subsets of J such that \(I' \cup I'' = I\) , \(\mathbf{I}' \cap \mathbf{I}'' = \varnothing, \mathbf{J}' \cup \mathbf{J}'' = \mathbf{J}, \mathbf{J}' \cap \mathbf{J}'' = \varnothing\) ; then there is a canonical associativity isomorphism (II, § 3, no. 9)

\[
\mathsf {T} _ {\mathbf {J}} ^ {\mathbf {I}} (\mathbf {M}) \rightarrow \mathsf {T} _ {\mathbf {J} ^ {\prime}} ^ {\mathbf {I} ^ {\prime}} (\mathbf {M}) \otimes_ {\mathbf {A}} \mathsf {T} _ {\mathbf {J} ^ {\prime \prime}} ^ {\mathbf {I} ^ {\prime \prime}} (\mathbf {M}).\tag{11}
\]

Considering the tensor algebra \(\mathsf{T}(\mathbf{M}\oplus \mathbf{M}^{*})\) , it follows from no. 5 that \(\mathsf{T}^n (\mathbf{M}\oplus \mathbf{M}^*)\) is canonically identified with the direct sum of the \(\mathsf{T}_{\mathbf{J}}^{\mathbf{I}}(\mathbf{M})\) where I runs through the set of subsets of the interval \([1,n]\) of \(\mathbf{N}\) and J is the complement of I in \([1,n]\) .

When \(I = [1, p]\) and \(J = [p + 1, p + q]\) with integers \(p \geqslant 0\) , \(q \geqslant 0\) (where we replace I (resp. J) by \(\varnothing\) when \(p = 0\) (resp. \(q = 0\) )), the A-module \(T_J^I(M)\) is also denoted by \(T_q^p(M)\) ; the A-modules \(T_0^n(M)\) and \(T_n^0(M)\) are therefore by definition the A-modules \(T^n(M)\) and \(T^n(M^*)\) respectively. When I and J are arbitrary finite sets of cardinals \(p = \text{Card}(I)\) and \(q = \text{Card}(J)\) , we give each of them a total ordering; then there exists an increasing bijection of I (resp. J) onto \([1, p]\) (resp. \([p + 1, p + q]\) ) and these bijections therefore define an isomorphism

\[
\mathbf {T} _ {\mathrm{J}} ^ {\mathrm{I}} (\mathbf {M}) \rightarrow \mathbf {T} _ {q} ^ {p} (\mathbf {M}).
\]

When M is a finitely generated projective A-module, it follows from II, § 2, no. 7, Corollary 4 to Proposition 13 and II, § 4, no. 4, Corollary 1 to Proposition 4 that there is a canonical isomorphism

\[
(\mathsf {T} _ {\mathbf {J}} ^ {\mathbf {I}} (\mathbf {M})) ^ {*} \to \mathsf {T} _ {\mathbf {I}} ^ {\mathbf {J}} (\mathbf {M}).
\]

Suppose now that \(\mathbf{M}\) is a finitely generated free A-module and let \((e_{\lambda})_{\lambda \in L}\) be a basis of \(\mathbf{M}\) (L therefore being a finite set). The basis of \(\mathbf{M}^*\) dual to \((e_{\lambda})\) (II, §2, no. 6) is denoted by \((e^{\lambda})_{\lambda \in L}\) . The bases \((e_{\lambda})\) and \((e^{\lambda})\) of \(\mathbf{M}\) and \(\mathbf{M}^*\) respectively define (no. 5) a basis of \(\mathsf{T}_{\mathbf{J}}^{\mathbf{I}}(\mathbf{M})\) , which we give explicitly as follows: given two mappings \(f\colon \mathbf{I} \to \mathbf{L}\) and \(g\colon \mathbf{J} \to \mathbf{L}\) , let \(e_f^g\) be the element \(\bigotimes_{i \in I \cup J} x_i\) of \(\mathsf{T}_{\mathbf{J}}^{\mathbf{I}}(\mathbf{M})\) defined by

\[
x _ {i} = e _ {f (i)} \quad \text { if } i \in \mathbf {I}, \qquad x _ {i} = e ^ {g (i)} \quad \text { if } i \in \mathbf {J}.
\]

When \((f, g)\) runs through the set of ordered pairs of mappings \(f: \mathbf{I} \to \mathbf{L}\) and \(g: \mathbf{J} \to \mathbf{L}\) , the \(e_f^g\) form a basis of the A-module \(\mathsf{T}_{\mathbf{J}}^{\mathbf{I}}(\mathbf{M})\) , said to be associated with the given basis \((e_{\lambda})\) of \(\mathbf{M}\) . For \(z \in \mathsf{T}_{\mathbf{J}}^{\mathbf{I}}(\mathbf{M})\) , we can therefore write

\[
z = \sum_ {(f, g)} \alpha_ {g} ^ {f} (z). e _ {f} ^ {g}\tag{12}
\]

where the \(\alpha_g^f\) are the coordinate forms relative to the basis \((e_f^g)\) ; by an abuse of language, the \(\alpha_g^f(z)\) are called the coordinates of the tensor \(z\) with respect to the basis \((e_\lambda)\) of the module M. The \(\alpha_g^f\) constitute the dual basis of \((e_g^f)\) , in other words they are identified with the elements of the basis of \(\mathsf{T}_{\mathbf{I}}^{\mathbf{j}}(\mathbf{M})\) , associated with \((e_\lambda)\) . When I and J are complementary subsets of an interval \([1,n]\) of N, \(\alpha_g^f\) (or \(\alpha_g^f(z)\) ) is denoted as follows: each element \(f(i)\) for \(i\in I\) is written as an upper index in the \(i\) -th place with a dot in the \(i\) -th place for \(i \in \mathbf{J}\) ; similarly, \(g(i)\) for \(i \in \mathbf{J}\) is written as a lower index in the \(i\) -th place with a dot in the \(i\) -th place for \(i \in \mathbf{I}\) . For example, for \(\mathbf{I} = \{1, 4\}\) , \(\mathbf{J} = \{2, 3\}\) , we write \(\alpha_{\cdot \nu \rho}^{\lambda \cdot \cdot \mu}\) if \(f(1) = \lambda, f(4) = \mu, g(2) = \nu, g(3) = \rho\) .

Let \((\bar{e}_{\lambda})_{\lambda \in L}\) be another basis of M and \(P\) the matrix of passing from \((e_{\lambda})\) to \((\bar{e}_{\lambda})\) (II, § 10, no. 8). Then the matrix of passing from \((e^{\lambda})\) to \((\bar{e}^{\lambda})\) (dual basis of \((\bar{e}_{\lambda}))\) is the contragredient \(^{t}P^{-1}\) of \(P\) (II, § 10, no. 8, Proposition 5). It follows (II, § 10, no. 10) that the matrix of passing from the basis \((e_{f}^{g})\) of \(\mathrm{T}_{\mathrm{J}}^{\mathrm{I}}(\mathrm{M})\) to the basis \((\bar{e}_{f}^{g})\) (where \(f\) (resp. \(g\) ) runs through the set of mappings of I into L (resp. of J into L)) is the matrix

\[
\bigotimes_ {i \in I \cup J} Q _ {i}, \quad \text { where } Q _ {i} = P \text { if } i \in I, Q _ {i} = ^ {t} P ^ {- 1} \text { if } i \in J.\tag{13}
\]

The transpose of this matrix is therefore identified with

\[
\bigotimes_ {i \in \mathbf {I} \cup J} R _ {i}, \quad \text { where } R _ {i} = ^ {t} P ^ {- 1} \text { if } i \in \mathbf {I}, R _ {i} = P \text { if } i \in \mathbf {J}.\tag{14}
\]

Suppose now that \(\mathbf{M}\) is an arbitrary module. Let \(i\in \mathbf{I}, j\in \mathbf{J}\) and write \(\mathbf{I}' = \mathbf{I} - \{i\}, J' = J - \{j\}\) ; we shall define a canonical A-linear mapping

\[
c _ {j} ^ {i} \colon \mathsf {T} _ {\mathbf {J}} ^ {\mathbf {I}} (\mathbf {M}) \to \mathsf {T} _ {\mathbf {J} ^ {\prime}} ^ {\mathbf {I} ^ {\prime}} (\mathbf {M}),
\]

called contraction of the index \(i\) and the index \(j\) . For this, note that the mapping of \(\mathbf{M}^{\mathrm{I}} \times (\mathbf{M}^{*})^{\mathrm{J}}\) , which associates with every family \((x_{i})_{i \in \mathbf{I} \cup \mathbf{J}}\) , where \(x_{i} \in \mathbf{M}\) if \(i \in \mathbf{I}\) and \(x_{i} \in \mathbf{M}^{*}\) if \(i \in \mathbf{J}\) , the element

\[
\left\langle x _ {i}, x _ {j} \right\rangle \otimes_ {k \in (\mathbb {I} \cup \mathbb {J}) - \{i, j \}} x _ {k}\tag{15}
\]

of \(T_{J'}^{I'}(M)\) , is A-multilinear; \(c_{j}^{i}\) is the corresponding A-linear mapping.

Suppose now that M is free and finitely generated and let \((e_{\lambda})_{\lambda \in L}\) be a basis of M. Given two mappings \(f: I \to L\) , \(g: J \to L\) , let \(f_i\) denote the restriction of f to \(I' = I - \{i\}\) and \(g_j\) the restriction of g to \(J' = J - \{j\}\) ; then by virtue of (12)

\[
c _ {j} ^ {i} (e _ {f} ^ {g}) = \left\{ \begin{array}{l l} 0 & \text { if } f (i) \neq g (j) \\ e _ {f _ {i}} ^ {g _ {j}} & \text { if } f (i) = g (j) \end{array} \right.\tag{16}
\]

The expression for the coordinates of \(c_j^i(z)\) as a function of those of \(z\) is obtained; for every mapping \(f'\) (resp. \(g'\) ) of \(I'\) into \(L\) (resp. of \(J'\) into \(L\) ) and every \(\lambda \in L\) , let \((f', \lambda)\) (resp. \((g', \lambda)\) ) denote the mapping of \(I\) into \(L\) (resp. of \(J\) into \(L\) ) whose restriction to \(I'\) (resp. \(J'\) ) is \(f'\) (resp. \(g'\) ) and which takes the value \(\lambda\) at the element \(i\) (resp. \(j\) ). Then, if the coordinate forms relative to the basis \((e_{f'}^{g'})\) of \(T_{J'}^{I'}(M)\) are denoted by \(\beta_{g'}^{f'}\) ,

\[
\beta_ {g ^ {\prime}} ^ {f ^ {\prime}} (c _ {j} ^ {i} (z)) = \sum_ {\lambda \in L} \alpha_ {(g ^ {\prime}, \lambda)} ^ {(f ^ {\prime}, \lambda)} (z).\tag{17}
\]

Examples of tensors. (1) Let M be a finitely generated projective A-module. We know (II, § 4, no. 2, Corollary to Proposition 2), that there is a canonical A-module isomorphism

\[
\theta_ {\mathbf {M}} \colon \mathbf {M} ^ {*} \otimes_ {\mathbf {A}} \mathbf {M} \rightarrow \operatorname{End} _ {\mathbf {A}} (\mathbf {M})
\]

such that \(\theta_{\mathbf{M}}(x^{*}\otimes x)\) (for \(x\in \mathbf{M},x^{*}\in \mathbf{M}^{*}\) ) is the endomorphism

\[
y \mapsto \langle y, x ^ {*} \rangle x.
\]

Hence, by means of \(\theta_{\mathbf{M}}\) , \(\mathsf{T}_{(1)}^{(2)}(\mathbf{M})\) (isomorphic to \(\mathsf{T}_1^1 (\mathbf{M}))\) can be identified with the A-module \(\operatorname{End}_{\mathbf{A}}(\mathbf{M})\) . Suppose that \(\mathbf{M}\) is a free module and let \((e_{\lambda})_{\lambda \in \mathbf{L}}\) be a basis of \(\mathbf{M}\) ; then the coordinates of a tensor \(z\in \mathbf{M}^*\otimes \mathbf{M}\) relative to the basis \((e^{\mu}\otimes e_{\lambda})\) of this module are denoted by \(\zeta_{\mu}^{\cdot \lambda}\) . As \(\theta_{\mathrm{M}}(e^{\mu}\otimes e_{\lambda})\) is the endomorphism \(y\mapsto \langle y,e^{\mu}\rangle e_{\lambda}\) , the endomorphism \(u = \theta_{\mathrm{M}}(z) = \theta_{\mathrm{M}}\left(\sum_{\lambda ,\mu}\zeta_{\mu}^{\cdot \lambda}e^{\mu}\otimes e_{\lambda}\right)\) maps \(y\) to \(\sum_{\lambda ,\mu}\zeta_{\mu}^{\cdot \lambda}\langle y,e^{\mu}\rangle e_{\lambda}\) ; writing \(y = e_{\lambda}\) , we obtain the relation

\[
u (e _ {\lambda}) = \sum_ {\mu \in L} \zeta_ {\lambda} ^ {\cdot \mu} e _ {\mu}\tag{18}
\]

in other words, the matrix of the linear mapping \(u = \theta_{\mathrm{M}}(z)\) is that whose element in the row of index \(\mu\) and column of index \(\lambda\) is \(\zeta_{\lambda}^{\mu}\) .

The definition of the trace of \(u\) (II, § 4, no. 3) shows immediately that

\[
\operatorname{Tr} (\theta_ {\mathbb {M}} (z)) = c _ {1} ^ {2} (z).
\]

Therefore the element \(z_0 = \sum_{\lambda \in L} e^\lambda \otimes e_\lambda\) (whose coordinates \(\zeta_\mu^{\cdot \lambda}\) are zero for \(\lambda \neq \mu\) and equal to 1 for \(\lambda = \mu\) ), which is such that \(\theta_M(z_0) = 1_M\) , is the image of the element \(1 \in A = T_0^0(M)\) under the mapping the transpose of the contraction

\[
c _ {1} ^ {2}: \mathbf {T} _ {\{1 \}} ^ {\{2 \}} (\mathbf {M}) \rightarrow \mathbf {A}.
\]

\begin{enumerate}
\def\labelenumi{(\arabic{enumi})}
\setcounter{enumi}{1}
\tightlist
\item
  Suppose always that M is a finitely generated projective A-module; there is a canonical A-module isomorphism
\end{enumerate}

\[
\mu : \mathbf {M} ^ {*} \otimes_ {\mathbf {A}} \mathbf {M} ^ {*} \rightarrow (\mathbf {M} \otimes_ {\mathbf {A}} \mathbf {M}) ^ {*}
\]

(II, § 4, no. 4, Corollary 1 to Proposition 2) and a canonical isomorphism

\[
\theta : (\mathbf {M} \otimes_ {\mathbf {A}} \mathbf {M}) ^ {*} \otimes_ {\mathbf {A}} \mathbf {M} \rightarrow \operatorname{Hom} _ {\mathbf {A}} (\mathbf {M} \otimes_ {\mathbf {A}} \mathbf {M}, \mathbf {M})
\]

(II, §4, no. 2, Corollary to Proposition 2); also \(\operatorname{Hom}_{\mathbf{A}}(\mathbf{M} \otimes_{\mathbf{A}} \mathbf{M}, \mathbf{M})\) is canonically isomorphic to the A-module \(\mathcal{L}_2(\mathbf{M}, \mathbf{M}; \mathbf{M})\) of A-bilinear mappings of \(\mathbf{M} \times \mathbf{M}\) into \(\mathbf{M}\) (II, §3, no. 9). Composing these isomorphisms, a canonical isomorphism is obtained

\[
\chi_ {\mathbf {M}} \colon T _ {\{1, 2 \}} ^ {(3)} (\mathbf {M}) = \mathbf {M} ^ {*} \otimes \mathbf {M} ^ {*} \otimes \mathbf {M} \rightarrow \mathcal {L} _ {2} (\mathbf {M}, \mathbf {M}; \mathbf {M})
\]

such that, for \(x^{*}, y^{*}\) in \(\mathbf{M}^{*}\) , \(z \in \mathbf{M}\) , \(\chi_{\mathbf{M}}(x^{*} \otimes y^{*} \otimes z)\) is the bilinear mapping

\[
(u, v) \mapsto \langle u, x ^ {*} \rangle \langle v, y ^ {*} \rangle z.
\]

Hence, by means of \(\chi_{\mathbf{M}}\) , \(\mathsf{T}_{\{1,2\}}^{(3)}(\mathbf{M})\) (isomorphic to \(\mathsf{T}_{2}^{1}(\mathbf{M})\) ) can be identified with the A-module \(\mathcal{L}_2(\mathbf{M},\mathbf{M};\mathbf{M})\) . Suppose that \(\mathbf{M}\) is a free A-module and let \((e_{\lambda})_{\lambda \in L}\) be a basis of \(\mathbf{M}\) ; then the coordinates of a tensor \(z\in \mathbf{M}^*\otimes \mathbf{M}^*\otimes \mathbf{M}\) relative to the basis \((e^{\lambda}\otimes e^{\mu}\otimes e_{\nu})\) of this module are denoted by \(\zeta_{\lambda \mu}^{\cdot \cdot \nu}\) . The bilinear mapping \(\chi_{\mathbf{M}}(z)\) maps the ordered pair \((e_{\lambda},e_{\mu})\) to

\[
\sum_ {\nu \in L} \zeta_ {\lambda \mu}^{\cdot \cdot \nu} e _ {\nu}
\]

and therefore the \(\zeta_{\lambda \mu}^{\cdots \nu}\) are just the constants of structure of the (in general non-associative) algebra defined on \(\mathbf{M}\) by the bilinear mapping \(\chi_{\mathbf{M}}(z)\) , with respect to the basis \((e_{\lambda})\) ( \(\S 1\) , no. 7).

Remark 2. Let \((e_{\lambda})_{\lambda \in L}, (\tilde{e}_{\lambda})_{\lambda \in L}\) be two bases of M and \(P\) the matrix of passage from \((e_{\lambda})\) to \((\tilde{e}_{\lambda})\) . On account of what was seen in Example 1, the element of \(P\) appearing in the row of index \(\lambda\) and the column of index \(\mu\) is denoted by \(\alpha_{\mu}^{\lambda}\) and the element of the contragredient \(^{t}P^{-1}\) appearing in the row of index \(\lambda\) and the column of index \(\mu\) is denoted by \(\beta_{\lambda}^{\mu}\) . Then (in the notation introduced above)

\begin{enumerate}
\def\labelenumi{(\arabic{enumi})}
\setcounter{enumi}{18}
\tightlist
\item
\end{enumerate}

\[
\left\{ \begin{array}{l} \tilde {e} _ {\mu} = \sum_ {\lambda} \alpha_ {\mu} ^ {\lambda} e _ {\lambda} \\ \tilde {e} ^ {\mu} = \sum_ {\lambda} \beta_ {\lambda} ^ {\mu} e ^ {\lambda} \end{array} \right.\tag{20}
\]

\[
\bar {e} _ {f ^ {\prime}} ^ {g ^ {\prime}} = \sum_ {(f, g)} \left(\prod_ {i \in I} \alpha_ {f ^ {\prime} (i)} ^ {f (i)}\right) \left(\prod_ {j \in J} \beta_ {g (j)} ^ {g ^ {\prime} (j)}\right) e _ {f} ^ {g}
\]

for all mappings \(f': \mathbf{I} \to \mathbf{L}\) and \(g': \mathbf{J} \to \mathbf{L}\) . The coordinates \(\zeta_g^f\) of a tensor \(z \in \mathbb{T}_J^{\mathbf{I}}(\mathbf{M})\) with respect to the basis \((e_{\lambda})\) can therefore be expressed in terms of the coordinates \(\bar{\zeta}_{g'}^{f'}\) of \(z\) with respect to the basis \((\bar{e}_{\lambda})\) using the formulae

\[
\zeta_ {g} ^ {f} = \sum_ {(f ^ {\prime}, g ^ {\prime})} \left(\prod_ {i \in I} \alpha_ {f ^ {\prime} (i)} ^ {f (i)}\right) \left(\prod_ {j \in J} \beta_ {g (j)} ^ {g ^ {\prime} (j)}\right) \bar {\zeta} _ {g ^ {\prime}} ^ {f ^ {\prime}}.\tag{21}
\]

The matrix \(P^{-1}\) of passage from the basis \((\tilde{e}_{\lambda})\) to the basis \((e_{\lambda})\) is the transpose of \(^{t}P^{-1}\) , so that \(\beta_{\lambda}^{\mu}\) is the element which appears in the column of index \(\lambda\) and the row of index \(\mu\) of \(P^{-1}\) . The calculation of \(e_{f}^{g}\) in terms of the \(\tilde{e}_{f'}^{g'}\) and that of the \(\overline{\zeta}_{g'}^{f'}\) in terms of the \(\zeta_{g}^{f}\) are therefore made by replacing \(\alpha_{\mu}^{\lambda}\) by \(\beta_{\lambda}^{\mu}\) and \(\beta_{\lambda}^{\mu}\) by \(\alpha_{\mu}^{\lambda}\) in the above calculations and exchanging the roles of \(f\) and \(f'\) and those of \(g\) and \(g'\) . Then

\[
e _ {f} ^ {g} = \sum_ {(f ^ {\prime}, g ^ {\prime})} \left(\prod_ {j \in J} \alpha_ {g ^ {\prime} (j)} ^ {g (j)}\right) \left(\prod_ {i \in I} \beta_ {f (i)} ^ {f ^ {\prime} (i)}\right) \bar {e} _ {f ^ {\prime}} ^ {g ^ {\prime}}\tag{22}
\]

\[
\bar {\zeta} _ {g ^ {\prime}} ^ {f ^ {\prime}} = \sum_ {(f, g)} \left(\prod_ {j \in J} \alpha_ {g ^ {\prime} (j)} ^ {g (j)}\right) \left(\prod_ {i \in I} \beta_ {f (i)} ^ {f ^ {\prime} (i)}\right) \zeta_ {g} ^ {f}.\tag{23}
\]

The above formulae are such that the summation is over indices which appear once as a lower index and once as an upper index. Certain authors allow themselves because of this circumstance to suppress the summation signs.

